%! Author = lili
%! Date = 25.09.25

To be able to contribute to the iGEM Wiki, it is beneficial for all team members to have a basic understanding  of \textbf{HTML\index{HTML}} and for at least the wiki team members to have a good understanding of \textbf{HTML\index{HTML}} and a basic understanding of \textbf{CSS \index{CSS}}.
Familiarity with these languages will enable everyone to assist with content creation, formatting, and layout
adjustments.

Consider organizing internal wiki workshops early in the season to teach the basics of HTML and CSS. This lowers the entry barrier for new members and enables more people to contribute to content creation and formatting. \\ \newline

\section{HTML\index{HTML}}\label{sec:html}
\paragraph{HTML\index{HTML}} is used to define the structure and content of a webpage.
It determines elements such as
headings, paragraphs, tables, images, links, and sections.
A clear and semantic HTML structure improves readability, maintainability, and accessibility. \\ \newline

There are many free resources available online, for example:
\begin{itemize}
    \item \url{https://www.w3schools.com/html/default.asp}
    \item \url{https://www.geeksforgeeks.org/html/html-cheat-sheet/}
\end{itemize}

If you are in a hurry, here are some key concepts to get started:

\begin{itemize} \index{Essentials}
\item \textbf{Opening and Closing Tags}: Each HTML\index{HTML} element typically consists of an opening tag and a closing tag.
For example, in \texttt{<p>Your text here</p>}, \texttt{<p>} is the opening tag, and \texttt{</p>} is the closing tag.

\item \textbf{Nesting of Tags}: HTML\index{HTML} tags can be nested within one another, meaning that one tag can contain another.
For instance, \texttt{<div><p>Your text here</p></div>}.

\item \textbf{Semantic HTML\index{HTML}}: Using appropriate tags for their intended purpose (like \texttt{<header>}, \texttt{<footer>}, \texttt{<article>}, etc.) improves accessibility and SEO.

\item \textbf{Attributes}: HTML\index{HTML} elements can have attributes that provide additional information about them.
For example, the \texttt{<a>} tag uses the \texttt{href} attribute to specify the link's destination: \texttt{<a href="https://example.com">Link</a>}.
\end{itemize}


\section{CSS\index{CSS} and SCSS\index{SCSS}}\label{sec:css-and-scss} \index{HTML!Styling}
\paragraph{CSS\index{CSS}} is used to style and arrange HTML elements.
It controls colors, fonts, spacing, layouts, animations, and responsive behavior for different screen sizes.
Even a basic understanding of CSS allows team members to make quick design adjustments and improve the overall appearance of the wiki. \\ \newline


\subsection*{Difference Between CSS and SCSS\index{SCSS}}

\textbf{SCSS\index{SCSS} (Sassy CSS\index{CSS})} is a syntax of SASS (Syntactically Awesome Style Sheets), which extends CSS with features like variables, nested rules, and mixins.
While CSS allows for styling web pages, SCSS\index{SCSS} offers more advanced capabilities for organizing and writing styles, making it easier to manage complex stylesheets.

\subsection*{How HTML\index{HTML} and CSS Work Together}

HTML\index{HTML} provides the structure and content of a webpage, while CSS handles the visual presentation.
HTML\index{HTML} elements are styled using CSS selectors, allowing developers to specify how different components should look.
For example, a \texttt{<h1>} element can be styled to change its font size, color, and alignment using CSS.


Resources:
\begin{itemize}
    \item \url{https://www.w3schools.com/css/default.asp}
    \item \url{https://www.geeksforgeeks.org/css/css-complete-guide/}
    \item \url{https://www.w3schools.in/scss/introduction}
\end{itemize}


Basic concepts for a quick start:
\begin{enumerate}
    \item \textbf{Selectors}: There are different types of CSS selectors (element, class, ID). For example, \texttt{.classname} targets all elements with that class, while \texttt{\#idname} targets a unique element.

    \item \textbf{Box model}: The box model describes how elements are rendered and how spacing around them is
    calculated, influencing overall layout.
    E.g.\ the aspects content, padding, border or margin

    \item \textbf{Flexbox and grid}: These modern CSS layout techniques help create responsive designs.

    \item \textbf{Responsive design through media queries}: Media queries enable styles to adapt based on screen size,
    orientation, or resolution, providing an optimal user experience across devices.
\end{enumerate}


\subsection*{Personal Variables}
Setting personal colors and variables in CSS is a good way to ensure design consistency, enhance maintainability, and make your styles easier to modify across an entire project. You can centralize design elements like colors, sizes, and gradients, making it simpler to update and reuse them throughout your code.
\\ \newline
In CSS, custom properties are defined using the \texttt{- -} prefix. For example, you might want to set up a collection of color variables that represent your project's brand or theme. These could include primary and secondary colors, as well as different shades for various purposes.
\begin{mybox}

    \begin{lstlisting}[language=XML]
:root {
  --text-primary: #850F78;
  --mediumpurple: #bc15aa;
  --lightpurple: #B85BD1;
  --very-light-purple: #ce9fc9;
}
    \end{lstlisting}
\end{mybox}
In this example, different shades of purple are defined as variables. The \texttt{:root} selector ensures that these variables are globally available throughout your CSS. Now, instead of repeating color codes in different parts of your stylesheet, you can reference these variables:
\begin{mybox}

    \begin{lstlisting}[language=XML]
h1 {
  color: var(--text-primary);
}

.button {
  background-color: var(--mediumpurple);
}

.highlight {
  border-color: var(--lightpurple);
}
    \end{lstlisting}
\end{mybox}
You can also use CSS variables to define more complex styles, such as gradients or sizes.
\begin{mybox}
    \begin{lstlisting}[language=XML]
:root {
  --node-size: 60px;
  --ourgradient: linear-gradient(90deg, rgba(147,94,184,1) 0%, rgba(160,167,243,1) 100%);
}
    \end{lstlisting}
\end{mybox}
Here, \texttt{--node-size} defines a consistent size that can be applied to elements like buttons, icons, or divs, while \texttt{--ourgradient} provides a linear gradient that can be used as a background for various components.


\begin{mybox}
    \begin{lstlisting}[language=XML]
.icon {
  width: var(--node-size);
  height: var(--node-size);
}

.header {
  background: var(--ourgradient);
}
    \end{lstlisting}
\end{mybox}

\section{JavaScript and TypeScript}\label{sec:javascript-and-typescript}

While HTML and CSS define the structure and appearance of a webpage, \textbf{JavaScript\index{JavaScript}} adds functionality and interactivity.
It allows websites to react to user input, update content dynamically, create animations, validate forms, or load data without reloading the page. \\ \newline

JavaScript can be used to create interactive diagrams, navigation menus, expandable content sections, image galleries, animated statistics, or custom user experiences.
Even small scripts can significantly improve usability and presentation quality. \\ \newline

\paragraph{TypeScript\index{TypeScript}} is an extended version of JavaScript that adds static typing and improved tooling.
It helps prevent programming errors, improves code readability, and becomes especially useful when several people collaborate on the same project.
For larger wiki projects, TypeScript is often preferable to plain JavaScript.
My template uses TypeScript\\ \newline

Relevant differences are:
\begin{itemize}
    \item \textbf{Typing}: JavaScript is dynamically typed, meaning variables can hold any data type, and type checking occurs at runtime.
    TypeScript is statically typed, allowing for type annotations and catching errors during development.

    \item \textbf{Development Experience}: TypeScript offers better tooling support, including autocompletion, type checking, and easier refactoring, leading to a more robust development experience.

    \item \textbf{Compilation}: TypeScript must be transpiled (compiled) to JavaScript before it can run in the browser, while JavaScript can be executed directly.

    \item \textbf{Language Features}: TypeScript supports modern JavaScript features and introduces additional functionalities not present in standard JavaScript, like interfaces and decorators.
\end{itemize}

%Wiki team members responsible for advanced features should understand variables, functions, events, DOM
%manipulation, modules, and debugging with browser developer tools. \\ \newline

If the team has little prior programming experience, it is perfectly reasonable to begin with JavaScript and later transition to TypeScript when the project grows in complexity. \\ \newline

Useful learning resources include:

\begin{itemize}
    \item \url{https://www.w3schools.com/typescript/index.php}
    \item \url{https://www.w3schools.com/js/default.asp}
    \item \url{https://www.geeksforgeeks.org/typescript/typescript-tutorial/}
    \item \url{https://www.geeksforgeeks.org/javascript/javascript-tutorial/}
\end{itemize}